\section{Introduction}
\label{sec:intro}

Test-Driven Development (TDD)~\cite{beck2002tdd} starts with a failing test, implements
behavior that passes it, and refactors while preserving the tests. Its usefulness depends
on whether those tests represent the intended behavior. An autonomous coding pipeline
can now perform the entire loop: a model translates a specification into tests, generates
code to pass them, and reviews both artifacts before reporting success. The pipeline
produces its own acceptance criteria and its own evidence of meeting them.

Tests and code can nevertheless share a misunderstanding. Practitioners describe tests
that reflect the implementation rather than the requirement as the \emph{ugly mirror},
with tautological tests as an extreme case~\cite{rudolph2008uglymirror,pereira2010ttdd}.
Writing tests first avoids copying an existing implementation but cannot guarantee a
correct interpretation of the task. We measure the \emph{verification gap}: tasks the
pipeline reports as solved whose implementations fail tests withheld from it
(Sec.~\ref{sec:gap}). This is an oracle-relative disagreement; case inspection is needed
to distinguish implementation defects from incomplete requirements or evaluation errors.

We instrument one test-first pipeline and vary which model writes and reviews its tests
and code. Starting with an inexpensive model in every role, we assign stronger models
to test authorship, review, or both while keeping the implementation generator fixed.
Experiments use selected HumanEval and LiveCodeBench problems and three frontier
verifiers. Every configuration retains false accepts. Some estimates favor stronger
verification, and the primary LiveCodeBench reviewer improves externally passing output,
but the false-accept comparisons do not establish a reliable reduction. Residual errors
do not establish that stronger verification cannot help.

RGRBench extends the study to Python packages with stateful and interacting components.
Here stronger review also rejects more work before code generation and reduces correct
output. Acceptance, correctness, and abandonment therefore need to be reported together.
Observed associations with task difficulty and retries do not establish causal effects.

Our contribution is an instrumented comparison of verification roles and pipeline verdicts,
including rejected and abandoned work, across these task settings. We release runners,
prompts, retained results, analyses, and a defect register, while disclosing protocol
deviations and limits of the reconstructed public history.\footnote{Artifact:
\url{https://github.com/integrallis/self-grading-gap}.
RGRBench: \url{https://github.com/integrallis/rgrbench}.}
